% Auto-generated complete chapter from 08_twinsight_x500_case_study.md
\chapter{08 — TwinSight X500 Case Study}
\label{complete:root_009}

\begin{sourcemeta}
\textbf{Fuente:} \href{file:///E:/Laboral/08_twinsight_x500_case_study.md}{08\_twinsight\_x500\_case\_study.md}\\
\textbf{Seccion:} root\\
\textbf{Estado:} canonical\\
\textbf{Rol:} Modulo principal\\
\textbf{Lineas originales:} 893\\
\textbf{Ultima modificacion:} 2026-06-11 13:32\\
\textbf{Deduplicacion conservadora:} 0 bloques repetidos omitidos; 0 caracteres repetidos no reimpresos.
\end{sourcemeta}

\section*{Resumen de orientacion}
This document defines the complete portfolio case study structure for \textbf{TwinSight X500 / WebGL-Thesis-Proposal}, the flagship project for Alexander Woodcock Salomón’s international job-search strategy.

\section*{Contenido completo depurado}
\addcontentsline{toc}{section}{Contenido completo depurado}




\subsection{Document purpose}




This document defines the complete portfolio case study structure for \textbf{TwinSight X500 / WebGL-Thesis-Proposal}, the flagship project for Alexander Woodcock Salomón’s international job-search strategy.




The objective is to convert the thesis project into a recruiter-readable and technically credible case study for these role families:




\begin{enumerate}
\item Real-Time 3D / Interactive Technical Visualization Developer.
\item Unity WebGL / Interactive 3D Developer.
\item Technical Visualization / Digital Twin / Simulation Developer.
\item Unity Technical Artist — optimization, runtime interaction, tools, technical UI.
\item CAD-to-Realtime Technical Artist / Optimization Specialist.
\end{enumerate}




This document is not the final public portfolio page yet. It is the \textbf{source case-study blueprint} from which the website page, GitHub README, LinkedIn project entry, ArtStation post, demo video, interview talking points and CV bullets should be derived.




\medskip\hrule\medskip




\subsection{Source-of-truth status}




\begin{center}
\scriptsize
\begin{longtable}{>{\raggedright\arraybackslash}p{0.455\linewidth} >{\raggedright\arraybackslash}p{0.455\linewidth}}
\toprule
Field & Current status \\
\midrule
\endfirsthead
\toprule
Field & Current status \\
\midrule
\endhead
Project name & TwinSight X500 / WebGL-Thesis-Proposal \\
Academic context & Ingeniería Multimedia thesis project, UNAD \\
Thesis status & Complete, pending defense / sustentación \\
Public repo & `\url{https://github.com/delarge95/WebGL-Thesis-Proposal`} \\
Candidate role & Unity integration, runtime systems, 3D optimization, UI, documentation, evaluation \\
Project type & Unity WebGL interactive technical visualization prototype \\
Subject & Holybro X500 V2 drone assembly visualization \\
Primary value & Convert complex 2D/CAD assembly documentation into interactive 3D understanding \\
Primary portfolio use & Flagship project for international technical art / real-time 3D / visualization roles \\
\bottomrule
\end{longtable}
\end{center}




\medskip\hrule\medskip




\subsection{Correct positioning}




\subsubsection{Recommended one-line positioning}




\textbf{TwinSight X500 is a Unity WebGL interactive technical visualization prototype that transforms complex drone assembly documentation into a real-time 3D inspection experience, combining CAD-to-realtime optimization, runtime interaction systems, technical UI and usability validation.}




\subsubsection{Short portfolio summary}




TwinSight X500 is an interactive 3D technical visualization system built in Unity WebGL for inspecting the assembly structure of a Holybro X500 V2 drone. The project converts manufacturer CAD and technical documentation into an optimized real-time model with component selection, exploded view, clipping, technical information panels and multiple visualization modes.




The project was developed as an academic thesis and validated with usability and workload instruments. Its purpose is not only to show parts, but to help users understand spatial relationships in a complex technical assembly.




\subsubsection{Recruiter-facing framing}




Use this project to signal:




\begin{itemize}
\item Unity / C\# implementation capacity.
\item Real-time 3D interaction design.
\item CAD-to-realtime asset optimization.
\item WebGL deployment awareness.
\item Technical UI and component information systems.
\item Ability to connect engineering documentation, 3D production and user evaluation.
\item Ability to document and defend technical design decisions.
\end{itemize}




\subsubsection{What this project should not claim}




Do \textbf{not} present TwinSight as:




\begin{itemize}
\item a commercial shipped product;
\item a production digital twin connected to live telemetry;
\item an engineering-grade simulation;
\item an FEA / CFD / thermodynamic simulator;
\item a certified aerospace training platform;
\item an original CAD design of the drone;
\item a full-stack SaaS product;
\item a solo-coded project without AI assistance;
\item evidence of senior-level industry production experience.
\end{itemize}




\medskip\hrule\medskip




\subsection{Core problem}




Traditional assembly documentation often relies on 2D diagrams, exploded drawings, text instructions and static CAD references. For spatially complex systems such as drones, users must mentally reconstruct relationships between parts, fasteners, modules, vertical layers and internal assemblies.




TwinSight addresses this problem by creating an interactive 3D representation where the user can inspect components, isolate assemblies, visualize relationships and understand how parts connect in space.




\subsubsection{Core thesis statement}




\textbf{Seeing parts is not enough. Users need to understand spatial relationships.}




\subsubsection{Problem dimensions}




\begin{enumerate}
\item \textbf{Spatial reconstruction load}  
\end{enumerate}

   Static documentation requires the user to infer 3D relationships from 2D views.




\begin{enumerate}
\item \textbf{Component relationship ambiguity}  
\end{enumerate}

   Assemblies contain parent parts, subparts, fasteners, mounts and nested modules that are difficult to understand from text alone.




\begin{enumerate}
\item \textbf{Scale and density}  
\end{enumerate}

   Imported CAD data can be too dense for real-time web deployment and must be optimized.




\begin{enumerate}
\item \textbf{Web performance constraints}  
\end{enumerate}

   A browser-based viewer must balance fidelity, interactivity and performance.




\begin{enumerate}
\item \textbf{Instructional clarity}  
\end{enumerate}

   Users need technical information, not just visual inspection.




\medskip\hrule\medskip




\subsection{Project objectives}




\subsubsection{General objective}




Design and implement an interactive Unity WebGL 3D visualization prototype that supports the inspection and understanding of the Holybro X500 V2 drone assembly.




\subsubsection{Specific objectives}




\begin{enumerate}
\item Convert CAD/manufacturer geometry and technical references into a real-time optimized 3D asset pipeline.
\item Implement runtime interaction systems for component selection, isolation, exploded view and technical inspection.
\item Design a technical UI that exposes part identity, specifications and assembly information.
\item Add visual modes that support different inspection needs: realistic, X-Ray, blueprint, solid color, wireframe, ghosted and qualitative thermal visualization.
\item Deploy the experience through Unity WebGL with web/mobile performance considerations.
\item Evaluate the prototype using usability, workload and qualitative feedback instruments.
\end{enumerate}




\medskip\hrule\medskip




\subsection{Project context and constraints}




\subsubsection{Technical context}




TwinSight is not a conventional game project. It sits between real-time 3D development, technical art, CAD optimization, digital product visualization and instructional interface design.




The central challenge is not only visual quality. The central challenge is making a complex technical assembly understandable, performant and navigable in a browser environment.




\subsubsection{Constraints}




\begin{center}
\scriptsize
\begin{longtable}{>{\raggedright\arraybackslash}p{0.455\linewidth} >{\raggedright\arraybackslash}p{0.455\linewidth}}
\toprule
Constraint & Impact \\
\midrule
\endfirsthead
\toprule
Constraint & Impact \\
\midrule
\endhead
CAD density & Required polygon reduction and selective reconstruction \\
WebGL target & Required performance-aware asset decisions \\
Academic scope & Prototype quality, not commercial deployment quality \\
Manufacturer assets & Requires careful attribution and portfolio risk mitigation \\
Time constraints & Some features may be prototype-level rather than production-hardened \\
AI-assisted coding & Requires transparent explanation of implementation process \\
\bottomrule
\end{longtable}
\end{center}




\medskip\hrule\medskip




\subsection{Technical stack}




\begin{center}
\scriptsize
\begin{longtable}{>{\raggedright\arraybackslash}p{0.455\linewidth} >{\raggedright\arraybackslash}p{0.455\linewidth}}
\toprule
Area & Tools / technologies \\
\midrule
\endfirsthead
\toprule
Area & Tools / technologies \\
\midrule
\endhead
Engine & Unity \\
Runtime language & C\# \\
Target platform & WebGL / browser \\
Rendering & URP, real-time visualization modes \\
UI & Unity UI Toolkit / technical interface systems \\
3D pipeline & Blender, CAD cleanup, retopology, optimization \\
CAD source & Holybro X500 V2 manufacturer geometry / technical references \\
Optimization & CAD-to-realtime reduction, modular assets, high-to-low workflow \\
Evaluation & SUS, NASA-TLX Raw, Think-Aloud, task records \\
Versioning & Git / GitHub \\
AI assistance & Used as implementation support; final integration/debugging/technical judgment retained by candidate \\
\bottomrule
\end{longtable}
\end{center}




\medskip\hrule\medskip




\subsection{Feature inventory}




\subsubsection{Core interaction features}




\begin{center}
\scriptsize
\begin{longtable}{>{\raggedright\arraybackslash}p{0.305\linewidth} >{\raggedright\arraybackslash}p{0.305\linewidth} >{\raggedright\arraybackslash}p{0.305\linewidth}}
\toprule
Feature & Status & Portfolio value \\
\midrule
\endfirsthead
\toprule
Feature & Status & Portfolio value \\
\midrule
\endhead
Component selection & Implemented / expected & Shows runtime interaction design \\
Technical part panels & Implemented / expected & Shows UI + documentation integration \\
Exploded view & Implemented / expected & Shows assembly relationship visualization \\
Cross-section / clipping & Implemented / expected & Shows technical inspection capability \\
Isolation / focus & Implemented / expected & Shows user-centered navigation \\
Multiple visual modes & Implemented / expected & Shows technical art and rendering modes \\
Mobile-first UI thinking & Implemented / expected & Shows web deployment awareness \\
Performance optimization & Implemented / expected & Shows CAD-to-realtime competence \\
\bottomrule
\end{longtable}
\end{center}




\subsubsection{Visual modes}




The public case study should show each visual mode with a screenshot or short clip:




\begin{enumerate}
\item Realistic.
\item X-Ray.
\item Blueprint.
\item Solid color.
\item Wireframe.
\item Ghosted.
\item Qualitative thermal visualization.
\end{enumerate}




\subsubsection{Technical data panels}




Recommended information hierarchy:




\begin{enumerate}
\item Identification.
\item Specifications.
\item Assembly relationship.
\item Parent / child relationship.
\item Notes / warnings / interpretation.
\end{enumerate}




\medskip\hrule\medskip




\subsection{Metrics and validation}




\subsubsection{Reported quantitative metrics}




\begin{center}
\scriptsize
\begin{longtable}{>{\raggedright\arraybackslash}p{0.305\linewidth} >{\raggedright\arraybackslash}p{0.305\linewidth} >{\raggedright\arraybackslash}p{0.305\linewidth}}
\toprule
Metric & Reported value & Interpretation \\
\midrule
\endfirsthead
\toprule
Metric & Reported value & Interpretation \\
\midrule
\endhead
CAD source complexity & Above 6.5M triangles & Indicates need for optimization \\
Optimized result & 95,617 triangles & Core evidence for CAD-to-realtime workflow \\
Validation participants & 12 & Formative usability validation; not large-scale statistical validation \\
SUS average & 91.88 & Strong usability signal, but sample size should be stated \\
NASA-TLX Raw — 3D viewer & 8.69 & Lower perceived workload condition \\
NASA-TLX Raw — 2D support & 19.89 & Higher perceived workload comparator \\
Task-condition records & 96 & Useful evaluation dataset size for task-level observations \\
\bottomrule
\end{longtable}
\end{center}




\subsubsection{How to present the metrics}




Use precise language:




\begin{quote}
In a formative academic evaluation with 12 participants and 96 task-condition records, the prototype obtained a SUS average of 91.88 and lower NASA-TLX Raw workload scores for the 3D viewer condition compared with the 2D support condition.
\end{quote}




Avoid overclaiming:




\begin{itemize}
\item Do not claim universal proof of superiority.
\item Do not claim statistical generalization unless the thesis includes formal statistical testing.
\item Do not claim commercial validation.
\item Do not claim training certification validity.
\end{itemize}




\subsubsection{Suggested visualization for portfolio}




Use a compact before/after metric panel:




\begin{center}
\scriptsize
\begin{longtable}{>{\raggedright\arraybackslash}p{0.455\linewidth} >{\raggedright\arraybackslash}p{0.455\linewidth}}
\toprule
Before / comparator & After / result \\
\midrule
\endfirsthead
\toprule
Before / comparator & After / result \\
\midrule
\endhead
Dense CAD above 6.5M triangles & Optimized model at 95,617 triangles \\
2D support NASA-TLX Raw 19.89 & 3D viewer NASA-TLX Raw 8.69 \\
Static assembly references & Interactive component inspection \\
\bottomrule
\end{longtable}
\end{center}




\medskip\hrule\medskip




\subsection{Candidate contribution map}




\subsubsection{Strongly defensible contribution areas}




\begin{center}
\scriptsize
\begin{longtable}{>{\raggedright\arraybackslash}p{0.455\linewidth} >{\raggedright\arraybackslash}p{0.455\linewidth}}
\toprule
Contribution & Evidence needed \\
\midrule
\endfirsthead
\toprule
Contribution & Evidence needed \\
\midrule
\endhead
Unity integration & repo structure, scripts, demo video \\
Runtime interaction systems & selection, exploded view, clipping, visual modes \\
Technical UI & screenshots, UI hierarchy, data panel structure \\
CAD-to-realtime workflow & before/after topology, polycount table, optimization notes \\
Blender optimization & screenshots of cleanup/retopo/high-low workflow \\
Propeller modeling & modeling breakdown / wireframe \\
Academic evaluation & SUS/TLX tables, methodology summary \\
Documentation & thesis excerpts, diagrams, README \\
AI-assisted implementation management & explanation of architecture, debugging, integration decisions \\
\bottomrule
\end{longtable}
\end{center}




\subsubsection{Carefully worded contribution statement}




\begin{quote}
I developed TwinSight X500 as a thesis prototype, leading the Unity integration, runtime interaction systems, technical UI, CAD-to-realtime optimization workflow, documentation and evaluation. AI tools were used as implementation support, but I was responsible for the architecture, integration, debugging, design decisions, validation and final project ownership.
\end{quote}




\subsubsection{Claims that require evidence before publication}




\begin{center}
\scriptsize
\begin{longtable}{>{\raggedright\arraybackslash}p{0.455\linewidth} >{\raggedright\arraybackslash}p{0.455\linewidth}}
\toprule
Claim & Evidence required \\
\midrule
\endfirsthead
\toprule
Claim & Evidence required \\
\midrule
\endhead
“Optimized from 6.5M+ to 95,617 triangles” & polycount screenshots / thesis table \\
“Unity WebGL deployment” & live build or recorded browser demo \\
“Mobile-first interface” & mobile screenshots / responsive demo \\
“Clipping system” & technical video or code snippet \\
“Visual modes” & screenshot grid \\
“SUS 91.88” & thesis evaluation section \\
“NASA-TLX reduction” & evaluation table \\
“AI-assisted but personally integrated” & interview explanation + code review readiness \\
\bottomrule
\end{longtable}
\end{center}




\medskip\hrule\medskip




\subsection{CAD-to-realtime pipeline narrative}




\subsubsection{Recommended pipeline section}




The case study should include a dedicated \textbf{CAD-to-Realtime Optimization Pipeline} section.




Suggested structure:




\begin{enumerate}
\item \textbf{Source}  
\end{enumerate}

   Manufacturer CAD / technical references from the Holybro X500 V2 assembly.




\begin{enumerate}
\item \textbf{Import / conversion}  
\end{enumerate}

   CAD imported or converted into polygonal geometry suitable for Blender and Unity.




\begin{enumerate}
\item \textbf{Assessment}  
\end{enumerate}

   Original geometry was too dense for WebGL deployment and contained parts requiring cleanup or reconstruction.




\begin{enumerate}
\item \textbf{Optimization}  
\end{enumerate}

   Reduced, rebuilt, retopologized or simplified geometry while preserving visual and instructional fidelity.




\begin{enumerate}
\item \textbf{Modularization}  
\end{enumerate}

   Organized the drone into inspectable component entities and scene nodes.




\begin{enumerate}
\item \textbf{Bake / material planning}  
\end{enumerate}

   High-to-low workflow used or planned for retaining detail on optimized meshes.




\begin{enumerate}
\item \textbf{Unity integration}  
\end{enumerate}

   Imported optimized assets into Unity, assigned metadata, interaction states and visual modes.




\begin{enumerate}
\item \textbf{Evaluation}  
\end{enumerate}

   Tested usability and cognitive workload in comparison with 2D support materials.




\subsubsection{Suggested asset breakdown}




\begin{center}
\scriptsize
\begin{longtable}{>{\raggedright\arraybackslash}p{0.455\linewidth} >{\raggedright\arraybackslash}p{0.455\linewidth}}
\toprule
Asset class & Portfolio treatment \\
\midrule
\endfirsthead
\toprule
Asset class & Portfolio treatment \\
\midrule
\endhead
Drone frame & show CAD-to-realtime conversion \\
Motors & show imported/manufacturer source + optimized version \\
Propellers & show original modeling contribution \\
Fasteners & show modular strategy, not every screw individually \\
Cables & show carefully if final model includes them; otherwise mark as scope limitation \\
Battery / GPS / telemetry modules & show if final and visually strong \\
\bottomrule
\end{longtable}
\end{center}




\medskip\hrule\medskip




\subsection{Runtime systems breakdown}




\subsubsection{System 1 — Component selection}




Public explanation:




\begin{quote}
The selection system allows users to inspect individual drone components and access part-specific technical information. Each selectable component is associated with metadata and visual feedback states.
\end{quote}




Evidence to show:




\begin{itemize}
\item selection highlight screenshot;
\item selected part panel;
\item hierarchy example;
\item code snippet only if clean and safe.
\end{itemize}




\subsubsection{System 2 — Exploded view}




Public explanation:




\begin{quote}
The exploded view reveals assembly relationships by separating components spatially while preserving their logical structure. This supports inspection of layered assemblies and internal connections.
\end{quote}




Evidence to show:




\begin{itemize}
\item normal vs exploded view comparison;
\item short animation clip;
\item note on limitations for cables/flexible elements if relevant.
\end{itemize}




\subsubsection{System 3 — Cross-section / clipping}




Public explanation:




\begin{quote}
The clipping system enables sectional inspection, allowing users to reveal internal structures and understand how components relate across vertical or horizontal planes.
\end{quote}




Evidence to show:




\begin{itemize}
\item clipping plane screenshot;
\item UI control screenshot;
\item simple diagram if available.
\end{itemize}




\subsubsection{System 4 — Visual modes}




Public explanation:




\begin{quote}
Multiple visual modes support different inspection needs: realistic rendering for recognition, ghosted/X-Ray for internal relationships, wireframe for structure and blueprint mode for technical interpretation.
\end{quote}




Evidence to show:




\begin{itemize}
\item visual mode grid;
\item explain purpose of each mode in one sentence.
\end{itemize}




\subsubsection{System 5 — Technical UI}




Public explanation:




\begin{quote}
The UI connects visual selection with structured component information. It is designed as an inspection interface rather than a game HUD.
\end{quote}




Evidence to show:




\begin{itemize}
\item panel layout;
\item part information hierarchy;
\item mobile/desktop variants if available.
\end{itemize}




\medskip\hrule\medskip




\subsection{Suggested public case study structure}




Use this order for the website / portfolio page.




\subsubsection{1. Hero section}




Content:




\begin{itemize}
\item Project title.
\item One-line summary.
\item 15–30 second autoplay muted video or high-quality render.
\item Tags: Unity, WebGL, C\#, Technical Visualization, CAD-to-Realtime, Blender, UI Toolkit.
\end{itemize}




Suggested hero copy:




\begin{quote}
Unity WebGL technical visualization prototype for inspecting the assembly structure of a Holybro X500 V2 drone.
\end{quote}




\subsubsection{2. Project snapshot}




Compact facts:




\begin{center}
\scriptsize
\begin{longtable}{>{\raggedright\arraybackslash}p{0.455\linewidth} >{\raggedright\arraybackslash}p{0.455\linewidth}}
\toprule
Category & Detail \\
\midrule
\endfirsthead
\toprule
Category & Detail \\
\midrule
\endhead
Type & Thesis / interactive technical visualization prototype \\
Role & Unity developer, technical artist, 3D optimization, UI, evaluation \\
Tools & Unity, C\#, WebGL, Blender, GitHub \\
Result & Optimized drone viewer with interaction, visual modes and usability validation \\
\bottomrule
\end{longtable}
\end{center}




\subsubsection{3. Problem}




Explain the documentation-to-understanding gap.




\subsubsection{4. Solution}




Explain the interactive viewer and main features.




\subsubsection{5. Technical pipeline}




Show CAD-to-realtime reduction and asset preparation.




\subsubsection{6. Runtime interaction systems}




Selection, exploded view, clipping, visual modes, UI.




\subsubsection{7. Evaluation results}




Show SUS, NASA-TLX Raw and task records.




\subsubsection{8. My contribution}




Explain what Alexander did, what came from manufacturer CAD and what was AI-assisted.




\subsubsection{9. Constraints and limitations}




Show maturity and honesty.




\subsubsection{10. What I would improve next}




Tie to employability and production mindset.




\medskip\hrule\medskip




\subsection{ArtStation / technical breakdown structure}




ArtStation should not be the main project home, but it can host a visual technical breakdown.




Recommended ArtStation post title:




\textbf{TwinSight X500 — CAD-to-Realtime Drone Visualization Breakdown}




Recommended sections:




\begin{enumerate}
\item Final viewer screenshots.
\item CAD source vs optimized model.
\item Wireframe / topology views.
\item Propeller modeling breakdown.
\item Material / visual mode grid.
\item Exploded view sequence.
\item Clipping / X-Ray inspection screenshots.
\item UI panels.
\item Performance/optimization summary.
\item Link to portfolio case study and GitHub.
\end{enumerate}




Avoid making ArtStation the only evidence source because recruiters for Unity/WebGL roles need technical explanation, not only renders.




\medskip\hrule\medskip




\subsection{GitHub README structure}




The GitHub README should be recruiter-readable and not thesis-only.




Recommended README sections:




\begin{enumerate}
\item Project name and one-line summary.
\item Live demo link or demo video link.
\item Screenshots.
\item Features.
\item Technical stack.
\item Architecture overview.
\item CAD-to-realtime optimization.
\item Evaluation metrics.
\item Installation / local run instructions.
\item Known limitations.
\item AI-assisted development disclosure.
\item Academic and asset attribution.
\item License / usage notes.
\end{enumerate}




\subsubsection{README warning}




If the repository contains messy commits, generated code, large assets, broken branches or academic clutter, do not overexpose it before cleanup. A polished README with curated screenshots is more important than making every internal artifact public.




\medskip\hrule\medskip




\subsection{Demo video plan}




\subsubsection{Recommended duration}




60–120 seconds.




\subsubsection{Target style}




Technical, direct, visual, no cinematic overproduction. The video should prove the system works.




\subsubsection{Structure}




\begin{center}
\scriptsize
\begin{longtable}{>{\raggedright\arraybackslash}p{0.455\linewidth} >{\raggedright\arraybackslash}p{0.455\linewidth}}
\toprule
Time & Content \\
\midrule
\endfirsthead
\toprule
Time & Content \\
\midrule
\endhead
0–10s & Final viewer hero shot; title overlay \\
10–25s & Problem: dense drone assembly / static documentation \\
25–45s & Component selection + technical panel \\
45–60s & Exploded view \\
60–75s & Clipping / cross-section \\
75–90s & Visual modes grid \\
90–105s & Optimization metric: 6.5M+ to 95,617 triangles \\
105–120s & Evaluation metrics + role statement \\
\bottomrule
\end{longtable}
\end{center}




\subsubsection{Suggested narration}




\begin{quote}
TwinSight X500 is a Unity WebGL technical visualization prototype developed for my Multimedia Engineering thesis. The goal was to transform complex drone assembly documentation into an interactive 3D inspection experience. I optimized manufacturer CAD data for real-time use, built runtime interaction systems for selection, exploded view and clipping, and designed technical UI panels for component-level information. In a formative evaluation with 12 participants, the prototype achieved a SUS average of 91.88 and lower perceived workload than the 2D support condition. The project demonstrates my focus on Unity, WebGL, CAD-to-realtime optimization and interactive technical visualization.
\end{quote}




\subsubsection{Demo video warnings}




Do not include:




\begin{itemize}
\item long thesis explanation;
\item raw editor footage unless useful;
\item unpolished UI states;
\item claims not visible in the video;
\item excessive AI/process discussion.
\end{itemize}




\medskip\hrule\medskip




\subsection{Screenshot and asset shot list}




\subsubsection{Mandatory screenshots}




\begin{enumerate}
\item Hero full drone view.
\item Selected component with technical panel.
\item Exploded view.
\item Cross-section / clipping.
\item X-Ray or ghosted mode.
\item Wireframe / blueprint mode.
\item Before/after CAD optimization if available.
\item UI close-up.
\item Mobile layout if available.
\item Evaluation metrics panel.
\end{enumerate}




\subsubsection{Optional screenshots}




\begin{enumerate}
\item Blender optimization workspace.
\item High-poly vs low-poly comparison.
\item UV layout.
\item Bake result.
\item Unity scene hierarchy.
\item Code architecture diagram.
\item Interaction flow diagram.
\end{enumerate}




\subsubsection{Avoid screenshots}




\begin{itemize}
\item messy Unity editor without explanation;
\item unfinished placeholder assets;
\item low-resolution mobile captures;
\item screenshots that reveal private notes, credentials or irrelevant folders;
\item copyrighted documentation pages unless clearly allowed and attributed.
\end{itemize}




\medskip\hrule\medskip




\subsection{Interview explanation}




\subsubsection{30-second answer}




\begin{quote}
TwinSight X500 is my thesis project: a Unity WebGL technical visualization prototype for inspecting a Holybro X500 V2 drone assembly. I focused on converting dense CAD/manufacturer geometry into a real-time optimized model, building interaction systems such as component selection, exploded view and clipping, and validating the result with usability and workload instruments. It is the clearest evidence of my fit for Unity, WebGL, technical art and interactive visualization roles.
\end{quote}




\subsubsection{90-second answer}




\begin{quote}
The project started from a documentation problem: complex assemblies are hard to understand through static 2D references alone. I built a Unity WebGL viewer that lets users inspect the Holybro X500 V2 drone in 3D, select components, view technical information, activate exploded view, use clipping and switch visual modes such as X-Ray, blueprint, wireframe and ghosted. The technical challenge was converting dense CAD-style geometry into a performant browser-ready experience while preserving the assembly meaning. I also evaluated the prototype with 12 participants using SUS, NASA-TLX Raw and Think-Aloud. The strongest result was that the 3D viewer condition showed lower perceived workload than the 2D support condition. AI tools helped with implementation support, but I handled architecture, integration, debugging, final design decisions and validation.
\end{quote}




\subsubsection{Answer to “How much did you build yourself?”}




\begin{quote}
I used AI tools as implementation support, mainly to accelerate coding, explore approaches and debug issues. I do not present the project as hand-coded line by line without assistance. My contribution was defining the architecture, integrating the Unity systems, adapting and debugging the code, building the interaction flow, preparing the 3D assets, connecting the UI and data structure, running the evaluation and making the final technical decisions. In interviews, I can explain the systems, trade-offs and limitations because I was responsible for the final integration and project outcome.
\end{quote}




\subsubsection{Answer to “Is this a commercial product?”}




\begin{quote}
No. It is an academic prototype developed as a thesis project. I treat it as a portfolio-quality technical visualization case study, not as a commercial training platform. The value is that it demonstrates an end-to-end workflow: CAD-to-realtime optimization, Unity WebGL implementation, technical UI and user evaluation.
\end{quote}




\subsubsection{Answer to “Is it a digital twin?”}




\begin{quote}
It is adjacent to digital twin visualization, but I would not call it a full digital twin because it does not include live data, bidirectional synchronization or operational telemetry. It is more accurately described as an interactive technical visualization prototype or assembly inspection viewer.
\end{quote}




\medskip\hrule\medskip




\subsection{AI-assisted development disclosure}




\subsubsection{Recommended public wording}




Use this only in GitHub or extended case study, not necessarily in the hero section:




\begin{quote}
AI tools were used as implementation support during development. Final architecture, integration, debugging, interaction design, asset preparation, evaluation and technical decisions were handled by the author.
\end{quote}




\subsubsection{When to disclose}




\begin{center}
\scriptsize
\begin{longtable}{>{\raggedright\arraybackslash}p{0.455\linewidth} >{\raggedright\arraybackslash}p{0.455\linewidth}}
\toprule
Context & Recommended level \\
\midrule
\endfirsthead
\toprule
Context & Recommended level \\
\midrule
\endhead
CV bullet & Usually omit unless asked \\
LinkedIn project & Usually omit or mention “AI-assisted tooling” only if relevant \\
GitHub README & Include short disclosure \\
Interview & Be transparent if asked \\
Portfolio case study & Include in “Process / Implementation” section, not hero \\
\bottomrule
\end{longtable}
\end{center}




\subsubsection{Risk control}




The disclosure must not sound like:




\begin{itemize}
\item “AI made the project.”
\item “I only prompted.”
\item “I cannot explain the code.”
\item “I copied generated systems without understanding.”
\end{itemize}




It must sound like:




\begin{itemize}
\item AI accelerated implementation.
\item Candidate owned architecture and integration.
\item Candidate can explain trade-offs.
\item Candidate validated the result.
\end{itemize}




\medskip\hrule\medskip




\subsection{Asset and IP risk mitigation}




\subsubsection{Risk}




The project uses manufacturer geometry or technical references related to the Holybro X500 V2. This is acceptable as academic/portfolio context only if handled carefully, but it requires attribution and avoidance of misleading ownership claims.




\subsubsection{Required public wording}




\begin{quote}
This project uses publicly available/manufacturer reference material for academic and portfolio purposes. The drone product and original manufacturer design belong to their respective owners. The portfolio case study focuses on the real-time visualization, optimization, interaction design and academic evaluation work.
\end{quote}




\subsubsection{Do not do}




\begin{itemize}
\item Do not imply ownership of the original Holybro drone design.
\item Do not redistribute proprietary CAD files if licensing is unclear.
\item Do not present manufacturer assets as original modeling work.
\item Do not include confidential documentation.
\item Do not use official logos in a way that implies endorsement.
\end{itemize}




\subsubsection{Safer publication strategy}




\begin{enumerate}
\item Publish optimized screenshots and demo video.
\item Publish code if clean and legally safe.
\item Avoid redistributing raw CAD sources.
\item Include attribution and academic-use note.
\item If uncertain, keep heavy assets out of the public repo.
\end{enumerate}




\medskip\hrule\medskip




\subsection{Role-specific positioning variants}




\subsubsection{Variant A — Real-Time 3D / Interactive 3D Developer}




Emphasis:




\begin{itemize}
\item Unity runtime systems.
\item WebGL deployment.
\item Interaction design.
\item Component selection.
\item Technical UI.
\end{itemize}




Suggested description:




\begin{quote}
Built an interactive Unity WebGL 3D inspection viewer for a drone assembly, implementing component selection, exploded view, clipping and visual modes for technical understanding.
\end{quote}




\subsubsection{Variant B — Unity Technical Artist}




Emphasis:




\begin{itemize}
\item CAD-to-realtime pipeline.
\item Optimization.
\item Visual modes.
\item Technical rendering states.
\item Bridge between assets and runtime.
\end{itemize}




Suggested description:




\begin{quote}
Developed a CAD-to-realtime optimization and visualization pipeline for a Unity WebGL drone assembly viewer, reducing dense source geometry into a browser-ready inspectable model.
\end{quote}




\subsubsection{Variant C — Technical Visualization / Digital Twin-adjacent}




Emphasis:




\begin{itemize}
\item Assembly inspection.
\item Component relationships.
\item Engineering documentation.
\item Evaluation.
\end{itemize}




Suggested description:




\begin{quote}
Created an interactive technical visualization prototype that helps users understand drone assembly relationships through 3D inspection, exploded view, clipping and structured component information.
\end{quote}




\subsubsection{Variant D — Unity WebGL Developer}




Emphasis:




\begin{itemize}
\item Browser deployment.
\item Performance constraints.
\item UI.
\item Runtime optimization.
\end{itemize}




Suggested description:




\begin{quote}
Implemented a Unity WebGL technical viewer designed for browser-based 3D inspection, balancing interaction, visual fidelity and optimized asset delivery.
\end{quote}




\subsubsection{Variant E — Game Technical Artist}




Emphasis:




\begin{itemize}
\item Technical art systems.
\item Visual modes.
\item Real-time rendering.
\item Optimization.
\end{itemize}




Warning: this should be secondary, because the project is not a game.




Suggested description:




\begin{quote}
Built real-time visualization systems in Unity, including X-Ray, wireframe, ghosted and blueprint modes, with an optimized 3D asset pipeline for interactive inspection.
\end{quote}




\medskip\hrule\medskip




\subsection{CV bullet bank}




Use only after the CV strategy section confirms the target role.




\subsubsection{Strong bullets}




\begin{itemize}
\item Developed \textbf{TwinSight X500}, a Unity WebGL technical visualization prototype for inspecting Holybro X500 V2 drone assembly relationships.
\item Optimized dense CAD/manufacturer geometry from \textbf{6.5M+ triangles to 95,617 triangles} for real-time browser-based inspection.
\item Implemented runtime interaction systems including component selection, exploded view, cross-section/clipping and multiple visualization modes.
\item Designed technical UI panels linking selected components to identification, specification and assembly information.
\item Conducted formative usability evaluation with \textbf{12 participants}, \textbf{96 task-condition records}, SUS and NASA-TLX Raw workload measures.
\item Achieved a \textbf{SUS average of 91.88} and lower perceived workload for the 3D viewer condition compared with 2D support material.
\item Built a CAD-to-realtime workflow using Blender and Unity to support optimized interactive technical visualization.
\end{itemize}




\subsubsection{Safer bullets if evidence is not yet public}




\begin{itemize}
\item Developed a Unity WebGL technical visualization thesis prototype for interactive drone assembly inspection.
\item Built component selection, exploded view, clipping and visual mode systems to support spatial understanding of a complex 3D assembly.
\item Prepared and optimized CAD-derived 3D assets for real-time visualization in a browser environment.
\item Evaluated the prototype through formative usability and workload testing.
\end{itemize}




\subsubsection{Avoid bullets}




\begin{itemize}
\item “Shipped a commercial digital twin platform.”
\item “Built an aerospace simulation system.”
\item “Created the original Holybro X500 V2 CAD model.”
\item “Developed a production-ready training platform.”
\item “Implemented AI-powered drone diagnostics.”
\end{itemize}




\medskip\hrule\medskip




\subsection{LinkedIn project description}




\subsubsection{Recommended version}




\textbf{TwinSight X500 — Unity WebGL Technical Visualization Prototype}




Developed a Unity WebGL interactive 3D visualization prototype for inspecting the assembly structure of a Holybro X500 V2 drone. The project converts dense CAD/manufacturer geometry and technical documentation into an optimized real-time viewer with component selection, exploded view, cross-section/clipping, technical information panels and multiple visualization modes.




The prototype was developed as a Multimedia Engineering thesis project and evaluated with usability and workload instruments including SUS, NASA-TLX Raw and Think-Aloud. The project demonstrates CAD-to-realtime optimization, Unity/C\# runtime systems, technical UI design and browser-based interactive 3D visualization.




\textbf{Technologies:} Unity, C\#, WebGL, URP, UI Toolkit, Blender, CAD-to-Realtime Optimization, Technical Visualization.




\medskip\hrule\medskip




\subsection{Portfolio page copy draft}




\subsubsection{Opening paragraph}




TwinSight X500 is a Unity WebGL technical visualization prototype designed to help users inspect and understand the assembly structure of a Holybro X500 V2 drone. The project addresses a common limitation of static technical documentation: users must mentally reconstruct complex 3D relationships from 2D diagrams, text and CAD references.




The prototype converts dense CAD/manufacturer geometry into an optimized real-time model and adds interaction systems for component selection, exploded view, clipping and visual mode switching. It was developed as a Multimedia Engineering thesis project and evaluated through a formative usability study.




\subsubsection{Problem paragraph}




Drone assemblies contain stacked frames, motors, fasteners, electronics, supports and internal relationships that are difficult to understand from static documentation alone. Traditional instructions show the parts, but they do not always make the spatial relationships clear. TwinSight focuses on making these relationships visible, inspectable and easier to reason about.




\subsubsection{Technical challenge paragraph}




The main technical challenge was converting dense CAD-style source geometry into a browser-ready interactive 3D experience. This required reducing geometric complexity, organizing components into inspectable entities, designing runtime interaction states and balancing visual clarity with WebGL performance constraints.




\subsubsection{Evaluation paragraph}




The prototype was evaluated with 12 participants using SUS, NASA-TLX Raw and Think-Aloud. The results indicated strong usability and lower perceived workload in the 3D viewer condition compared with the 2D support condition. These findings should be interpreted as formative academic validation rather than commercial-scale product validation.




\medskip\hrule\medskip




\subsection{Evidence checklist before publication}




\subsubsection{Minimum publication-ready evidence}




\begin{center}
\scriptsize
\begin{longtable}{>{\raggedright\arraybackslash}p{0.455\linewidth} >{\raggedright\arraybackslash}p{0.455\linewidth}}
\toprule
Item & Required before public launch \\
\midrule
\endfirsthead
\toprule
Item & Required before public launch \\
\midrule
\endhead
60–120 sec demo video & Yes \\
Hero screenshot & Yes \\
Feature screenshot grid & Yes \\
Optimization before/after metric & Yes \\
SUS / NASA-TLX summary & Yes \\
Attribution note & Yes \\
GitHub README cleanup & Yes \\
AI-assisted disclosure & Yes, at least in GitHub/process \\
Broken branches removed/hidden & Yes \\
Live WebGL build & Strongly recommended \\
\bottomrule
\end{longtable}
\end{center}




\subsubsection{Optional but high-value evidence}




\begin{center}
\scriptsize
\begin{longtable}{>{\raggedright\arraybackslash}p{0.455\linewidth} >{\raggedright\arraybackslash}p{0.455\linewidth}}
\toprule
Item & Value \\
\midrule
\endfirsthead
\toprule
Item & Value \\
\midrule
\endhead
Live browser demo & Very high \\
Interaction architecture diagram & High \\
CAD-to-realtime pipeline diagram & High \\
Technical UI flow diagram & Medium-high \\
Code snippet of clean system & Medium \\
Mobile performance notes & Medium \\
Thesis PDF excerpt & Medium, if clean \\
\bottomrule
\end{longtable}
\end{center}




\medskip\hrule\medskip




\subsection{Known limitations section}




The public case study should include limitations to avoid overclaiming.




Recommended wording:




\begin{quote}
TwinSight X500 is an academic prototype, not a commercial training product. It does not include live telemetry, physics simulation, certified maintenance workflows or production-grade backend infrastructure. The project focuses on interactive assembly visualization, CAD-to-realtime optimization, technical UI and formative usability evaluation.
\end{quote}




Potential limitations:




\begin{itemize}
\item WebGL performance may vary by device.
\item Some geometry may be simplified for real-time use.
\item Some flexible elements such as cables may be simplified or excluded.
\item Evaluation sample size is formative.
\item Manufacturer assets require careful attribution.
\item The system is not a full digital twin.
\end{itemize}




\medskip\hrule\medskip




\subsection{Final case study decision}




TwinSight X500 should be the \textbf{primary flagship project} across:




\begin{enumerate}
\item Portfolio homepage.
\item GitHub pinned repositories.
\item LinkedIn Featured section.
\item CV project section.
\item Demo reel.
\item Interview narrative.
\end{enumerate}




It is the project with the strongest alignment between Alexander’s real profile and the target labor market: Unity, WebGL, CAD-to-realtime, real-time interaction, technical visualization, UI, Blender optimization and academic evaluation.




ARA, Blender portrait and other projects should support the profile, but they should not compete with TwinSight for attention in the first 10 seconds of recruiter review.




\medskip\hrule\medskip




\subsection{Completion status}




\begin{center}
\scriptsize
\begin{longtable}{>{\raggedright\arraybackslash}p{0.455\linewidth} >{\raggedright\arraybackslash}p{0.455\linewidth}}
\toprule
Section & Status \\
\midrule
\endfirsthead
\toprule
Section & Status \\
\midrule
\endhead
Positioning & Complete \\
Problem framing & Complete \\
Technical stack & Complete \\
Feature inventory & Complete \\
Metrics framing & Complete \\
Contribution map & Complete \\
Demo plan & Complete \\
Screenshot plan & Complete \\
GitHub README outline & Complete \\
LinkedIn project draft & Complete \\
Risk mitigation & Complete \\
Evidence checklist & Complete \\
\bottomrule
\end{longtable}
\end{center}




\medskip\hrule\medskip




\subsection{Next actions for this project}




\begin{enumerate}
\item Record a 60–120 second demo video.
\item Capture the mandatory screenshot set.
\item Produce a visual mode grid.
\item Capture CAD-to-realtime before/after evidence.
\item Clean the GitHub README.
\item Add a public attribution and AI-assisted implementation disclosure.
\item Build the portfolio page using the case study structure above.
\item Prepare the 30-second and 90-second interview explanations.
\end{enumerate}



